\properlane{purified-desire}{Purified desire and the sight of God}
\subsection{A gift already received and a promise still awaited}
John's short passage makes the believer both a recipient and an agent. The Father's love makes children of God. Those children then purify themselves because they hope to see him. The apparent difficulty is sharper than the familiar observation that people need to improve. If God gives the identity, why does its fulfillment still require purification? And if the believer purifies himself, what remains of the claim that the Father gives everything?

Augustine begins with the reality of sonship. A name without the thing named would be useless, like calling someone a physician who cannot heal. God's children must be such in their life, faith, hope and charity. Yet Augustine does not infer that their present condition exhausts what John promises. The phrase ``not yet'' describes a future manifestation of the life they already have. Their sonship has begun truly while the promised sight exceeds their present experience.\footnote{Augustine, \work{Tractates on First John} IV.4--5.}

His distinction between two ways of seeing Christ explains the difference. Christ's humanity is visible, and those condemned at judgment will see the human judge. John's promise cannot mean merely that encounter, since he connects the promised vision to likeness and purification. Augustine consequently speaks of seeing Christ as the divine Word, equal with the Father. The difference concerns what is seen and the beholder's relation to him. A visual encounter with a human figure would not by itself account for the transforming likeness John announces.

The psalm asks who can ascend and stand in God's holy place. Its question can be read beside this promise: approach has a destination, and that destination is communion with God. Aquinas's distinction between ascent and remaining makes this precise. An intention to seek him begins a journey; standing concerns its completion. In this interpretation, the saints celebrated at the Entrance are those whose journey gives the congregation reason to rejoice, while the congregation's Collect still asks for mercy along its own way. That connection follows from the difference between present adoption and completed sight.\footnote{Aquinas, \work{Super Psalmos} 23, n. 2, paragraph 86989; 1 John 3:1--2.}

\subsection{Why desire must grow}
Augustine explains waiting by an ordinary container. A bag or skin intended to receive something large must be stretched. The promised gift does not change; the capacity of its recipient grows. He then changes the comparison from size to contents. A vessel full of vinegar has to be emptied and cleaned before honey can be poured in. The two comparisons answer different parts of the problem. Desire must expand, and desires that exclude the promised good must be relinquished.\footnote{Augustine, \work{First John} IV.6.}

His examples concern a good that words cannot adequately describe. Honey, gold and wine are comparisons, and he explicitly says they fall short of God. The purpose of desire is consequently more than increasing the intensity of an appetite for a familiar object. The person learns to seek the one from whom created beauty comes. Waiting can expose how readily a person identifies happiness with something smaller and already possessed.

The psalm's refusal of vanity has a similar moral consequence. Augustine understands the soul as destined for an enduring life and therefore desiring an unchangeable good. In the psalm, clean hands keep this desire joined to deeds. A person cannot seek an eternal good only in imagination while treating neighbors unjustly. Aquinas likewise distinguishes purity from disordered interior appetites and warns, in his explanation of the verse, about pride arising even from an apparent purity. A purified heart can still be compromised by taking pride in its own condition.\footnote{Augustine, \work{Expositions}, Psalm 23.4; Aquinas, \work{Super Psalmos} 23, n. 3, paragraph 86990.}

The Prayer over the Offerings concerns the saints' solicitude for people still seeking salvation. Beside Augustine's account, that intercession means that completed desire for God is compatible with concern for others. A person perfected in charity does not become less capable of loving those who need help. The prayer's subject therefore corrects an account of contemplation that would end in isolation. The proper Preface's help and example similarly concern a pilgrimage undertaken in communion.

\subsection{Who performs the purification?}
Augustine directly asks who purifies the believer. His answer is God. He then explains why John still attributes purification to the person. God does not purify an unwilling subject: the believer joins his will to God's action. Consent is a real act, although the indwelling helper makes the purification possible. Augustine uses the request for help to show both sides. Asking for a helper acknowledges need, and the word ``help'' implies that the person asking also acts.\footnote{Augustine, \work{First John} IV.7.}

This account gives the Collect a clear relation to the epistle. Asking for mercy is an act of the needy will directed toward God. The saints' intercession does not remove the assembly's own petition or its need to consent. Likewise the Lamb's blood in Revelation identifies the saving cause of the robes' whiteness. Bede's explanation extends the cleansing to the whole Church. The shared inference is that holiness received from Christ can require the recipient's fidelity without making that person the independent source of holiness.\footnote{Bede, \work{Explanation}, I, on 7:14, p. 53.}

Aquinas reaches a compatible result through another text. In Psalm 23:5, blessing and mercy ordinarily appear after the description of innocence. Their position might suggest only a payment for prior achievement. Aquinas also permits the causal explanation: the person approaches because God's blessing and mercy make innocence possible. The same verses therefore allow him to explain both reward and enabling grace. This is a separate exegetical argument from Augustine's discussion of the will, and it answers the psalm's own problem of access.\footnote{Aquinas, \work{Super Psalmos} 23, n. 4, paragraph 86991.}

The acclamation's invitation to the burdened belongs here because Christ calls people to receive rest from him. Chrysostom explains the yoke in the following context through learning humility and meekness. The willing person still labors, but learns from Christ rather than presuming to achieve the end by self-sufficiency. The invitation precedes the demand in Matthew 11:28--30. When heard before the Beatitudes, it allows their difficult demands to be understood within Christ's call and assistance. Chrysostom himself relates that humility to the beginning of Jesus' teaching.\footnote{Chrysostom, \work{Matthew} XXXVIII.3.}

\subsection{Why the pure heart is promised vision}
Aquinas develops the difficulty in Matthew 5:8 by setting the promise beside texts that call God unseen and inaccessible. One solution would deny the sight of God's essence and offer only a created reflection of his brightness. Aquinas rejects it because John promises seeing him as he is. He also argues from desire: someone who knows effects seeks their cause, and a final happiness that left the first cause unknown would leave that desire unfinished.\footnote{Aquinas, \work{Super Matthaeum}, cap. 5, lectio 2, paragraph 87340, on Matthew 5:8, citing 1 John 3:2 and 1 Corinthians 13:12.}

A second solution would make the divine essence visible to bodily eyes. Aquinas rejects that explanation through the Beatitude's own wording. Jesus names the heart. The mode of seeing is intellectual, as the scriptural language of the eyes of the heart suggests. An increase in a bodily eye's power would not turn the divine essence into a colored object. Aquinas can still speak of glorified creation manifesting God with extraordinary clarity, while distinguishing that manifestation from the intellectual sight of God himself.

These arguments clarify why purity concerns the heart and why the promise is future. Aquinas distinguishes acts of virtue that dispose a person toward happiness from the happiness they seek. Moral action orders the person toward purity and peace. Completed happiness involves the vision of God and the union of love. He therefore cites both John's promise of sight and John's statement about the Father's love and sonship. Knowledge and loving communion occur together in his account.\footnote{Aquinas, \work{Super Matthaeum}, cap. 5, lectio 2, paragraph 87340, opening treatment of the Beatitudes and comments on 5:8--9.}

This distinction does not empty present life of blessedness. Aquinas describes beginnings now and consummation in the heavenly homeland. Augustine likewise understands the promises as truly begun in the saints' present life. The future is fuller than the present without making present grace fictitious. The Prayer after Communion can consequently ask for charity's completion after the assembly has truly received the sacrament. Its hoped-for heavenly banquet concerns the fulfillment of communion, and does not make present reception an automatic assurance of perseverance.

\subsection{The Beatitudes as a course of conversion}
Augustine begins the course with humility because pride resists God's teaching. Meekness then concerns a teachable reception of Scripture. Mourning follows the discovery of one's entanglement in lesser loves. Hunger for justice concerns the difficult effort of leaving those attachments. Mercy appears as the fitting conduct of someone who needs superior help and can still help someone weaker. Purity and peace concern a person whose desires are being ordered toward God.\footnote{Augustine, \work{Sermon on the Mount} I.3--4, §§10--12. His order places meekness before mourning, as does the identified Douay--Rheims study witness; the U.S. proclaimed Gospel places mourning first.}

The explanation of mercy is particularly concrete. A person asking help from God may be stronger than a neighbor in some matter. Augustine reasons that the one seeking help should give the help he can. Thus dependence on grace gives a reason for mercy instead of excusing inaction. His course also explains why purity cannot be detached from the preceding works of charity. The heart learns singleness through a conversion that includes conduct toward others.

His account of peace concerns the ordering of desires under reason and of reason under God. It describes the person in whom disordered appetite no longer rules. Persecution then reveals whether this interior order can endure an attack from outside. The promise of the kingdom at the first and eighth blessings lets Augustine relate the conclusion to the beginning. The eighth tests the life developed through the preceding blessings.

At Communion, the selection of purity, peace and persecution permits that course to be considered at its later stages. The congregation receives sacramentally while still seeking the completed vision; it seeks sonship's likeness while already adopted. If the solemn blessing is chosen, its concern for example and service extends that hope into conduct. The earthly pilgrimage described in the Preface requires the active charity for which the final prayer asks.

\subsection{The four senses}
\begin{fourSenses}
\item[Literal.] John addresses actual children of God who await manifestation and purify themselves in hope. The psalm asks about admission to God's holy place; Jesus promises the sight of God to the clean of heart; Revelation describes a multitude cleansed through the Lamb. Their different literary settings remain distinct.
\item[Allegorical.] Christ's saving blood makes the Church the people seeking the Father's face. Augustine reads the psalm through Christ's resurrection and glorification. Aquinas brings Matthew's promised vision together with John's promise of likeness, as the life of those united to the Son.
\item[Moral.] The will consents to God's purifying help. Clean hands, undivided desire, mercy toward those needing assistance and endurance belong to that consent. Augustine's account of desire explains why relinquishing lesser attachments concerns preparation for communion with God.
\item[Anagogical.] The promised fulfillment is seeing God and sharing the heavenly banquet in perfected charity. The present beginnings of blessedness remain incomplete. Augustine distinguishes creaturely likeness from equality with God, even when that likeness is perfected.\footnote{Augustine, \work{First John} IV.9.}
\end{fourSenses}
